\documentclass[11pt]{article}
\usepackage[T1]{fontenc}
\usepackage[letterpaper,margin=0.9in]{geometry}
\usepackage{amsmath,amssymb,amsthm}
\usepackage{graphicx}
\usepackage[hidelinks]{hyperref}
\setlength{\parindent}{0pt}
\setlength{\parskip}{0.6em}

\newtheorem{theorem}{Theorem}
\newtheorem{corollary}[theorem]{Corollary}
\newtheorem{remark}[theorem]{Remark}
\newcommand{\D}{\Delta}

\title{A product theorem for $\Delta$-spaces\\with a point-finite-assignment factor}
\author{Candidate partial result for arXiv:2011.14299}
\date{11 August 2026}

\begin{document}
\maketitle

\textbf{Status:} \texttt{candidate\_partial\_result\_likely\_valid}.

\section{Source problem and later status}

K\k{a}kol and Leiderman ask whether the product of two compact
$\Delta$-spaces must again be a $\Delta$-space
\cite[Problem~5.9, PDF p.~12]{KL}.  (Later references use Problem~5.10,
reflecting a numbering change in the journal version.)

\begin{center}
\includegraphics[width=0.96\textwidth]{figures/open_problem_crop.png}
\end{center}

Leiderman and Szeptycki explicitly report that the question remains open even
for compact factors \cite[Section~5]{LS}.  The result below gives a positive
answer whenever one factor has a point-finite neighborhood assignment, and a
countable-closed-cover extension reaches known non-Eberlein compact examples.

\section{Definitions and idea of the proof}

A Tychonoff space $X$ is a \emph{$\Delta$-space} if every decreasing sequence
of subsets with empty intersection admits a decreasing open expansion with
empty intersection.  We use the equivalent characterization from
\cite[Theorem~2.1]{KL}: every countable pairwise disjoint family of subsets of
$X$ admits a point-finite open expansion.

A \emph{point-finite neighborhood assignment} (PFNA) on $Y$ is a family
$\{N_y:y\in Y\}$ such that $N_y$ is an open neighborhood of $y$ and, for each
$t\in Y$, only finitely many $y$ satisfy $t\in N_y$.

\subsection*{Idea of the proof}

Given disjoint sets in $X\times Y$, freeze $y\in Y$ and expand their vertical
sections point-finitely in $X$.  Thicken each fiberwise expansion by the
assigned neighborhood $N_y$ and take the union over $y$.  At a product point
$(x,t)$, point-finiteness on $Y$ leaves only finitely many relevant fibers;
point-finiteness in $X$ leaves only finitely many indices on each such fiber.
This two-stage finiteness is exactly what the product problem normally lacks.

\section{The fiberwise product theorem}

\begin{theorem}[PFNA factor theorem]\label{thm:pfna}
Let $X$ be a $\Delta$-space and let $Y$ be a Tychonoff space admitting a
point-finite neighborhood assignment.  Then $X\times Y$ is a $\Delta$-space.
\end{theorem}

\begin{proof}
Fix a pairwise disjoint sequence $\{A_n:n\in\mathbb N\}$ of subsets of
$X\times Y$.  For $y\in Y$, let
\[
 A_n(y)=\{x\in X:(x,y)\in A_n\}.
\]
For fixed $y$, the sets $A_n(y)$ are pairwise disjoint.  Since $X$ is a
$\Delta$-space, choose a point-finite family of open sets
$\{U_{n,y}:n\in\mathbb N\}$ in $X$ such that
$A_n(y)\subseteq U_{n,y}$ for every $n$.

Let $\{N_y:y\in Y\}$ be a PFNA and define
\[
 W_n=\bigcup_{y\in Y} U_{n,y}\times N_y.
\]
Each $W_n$ is open.  If $(x,y)\in A_n$, then $x\in U_{n,y}$ and
$y\in N_y$, so $(x,y)\in W_n$.

It remains to check point-finiteness.  Fix $(x,t)\in X\times Y$ and put
\[
 F_t=\{y\in Y:t\in N_y\}.
\]
The set $F_t$ is finite.  Moreover,
\[
 \{n:(x,t)\in W_n\}
 \subseteq
 \bigcup_{y\in F_t}\{n:x\in U_{n,y}\}.
\]
Every set on the right is finite by the chosen fiberwise expansion, and there
are only finitely many of them.  Thus $\{W_n:n\in\mathbb N\}$ is point-finite.
The disjoint-family characterization now gives $X\times Y\in\Delta$.
\end{proof}

\begin{corollary}\label{cor:closedcover}
Suppose $X\in\Delta$ and
$Y=\bigcup_{m\in\mathbb N}Y_m$, where every $Y_m$ is closed in $Y$ and admits
a PFNA.  Then $X\times Y\in\Delta$.
\end{corollary}

\begin{proof}
Theorem~\ref{thm:pfna} gives $X\times Y_m\in\Delta$ for every $m$, and these
subspaces are closed in $X\times Y$.  K\k{a}kol and Leiderman's countable
closed-union theorem \cite[Proposition~2.2]{KL2} now yields the conclusion.
\end{proof}

\section{Compact consequences}

Every scattered Eberlein compactum admits a PFNA
\cite[Remark~2.4 and the discussion after Theorem~3.7]{KL}; see also
\cite[Section~6]{LS}.  Therefore:

\begin{corollary}\label{cor:compact}
Let $X$ be any $\Delta$-space.
\begin{enumerate}
\item If $K$ is a scattered Eberlein compact space, then $X\times K\in\Delta$.
\item If a compact space $K$ is a countable union of scattered Eberlein
compact subspaces, then $X\times K\in\Delta$.
\item In particular, the conclusion holds for every compact scattered $K$ of
Cantor--Bendixson height at most two.
\end{enumerate}
\end{corollary}

\begin{proof}
Part (1) is Theorem~\ref{thm:pfna}.  In (2), the compact subspaces are closed,
so Corollary~\ref{cor:closedcover} applies.  A compact scattered space of
height at most two is a finite topological sum of one-point compactifications
of discrete spaces, up to finitely many isolated points; hence it is a
scattered Eberlein compactum.
\end{proof}

The next case is useful because the factor need not itself be Eberlein.  Let
$\mathcal A$ be an almost disjoint family of infinite subsets of $\mathbb N$,
let $\Psi(\mathcal A)=\mathbb N\cup\mathcal A$ be the corresponding
Isbell--Mr\'owka space, and let $K(\mathcal A)$ be its one-point
compactification.

\begin{corollary}\label{cor:psi}
For every $\Delta$-space $X$ and every almost disjoint family
$\mathcal A\subseteq[\mathbb N]^{\omega}$,
\[
 X\times K(\mathcal A)\in\Delta.
\]
When $\mathcal A$ is uncountable, this includes the non-Eberlein compacta used
in \cite[Theorem~3.10]{KL}.
\end{corollary}

\begin{proof}
Inside $K(\mathcal A)$, the set
$K_0=\mathcal A\cup\{\infty\}$ is closed and is the one-point
compactification of the discrete space $\mathcal A$.  Thus $K_0$ is a
scattered Eberlein compactum and admits a PFNA.  The remaining points are the
countably many closed singletons $\{n\}$, $n\in\mathbb N$.  Hence
\[
 K(\mathcal A)=K_0\cup\bigcup_{n\in\mathbb N}\{n\}
\]
is a countable closed PFNA cover, and Corollary~\ref{cor:closedcover} applies.
\end{proof}

The same proof works for a generalized $\Psi(D,\mathcal A)$ whenever
$D=\bigcup_mD_m$ and every $D_m$ is closed discrete in
$\Psi(D,\mathcal A)$: use the closed compact PFNA pieces
$\mathcal A\cup\{\infty\}$ and $D_m\cup\{\infty\}$.

\section{Upgrade audit, limits, and novelty}

Eight focused upgrade routes were checked.  The successful chain was:
fiberwise PFNA factor theorem; scattered Eberlein specialization; countable
closed PFNA covers; height-two compacta; compactified classical
$\Psi$-spaces; and the generalized $\Psi$ exhaustion above.  Two routes did
not close the full problem.  First, proving that every compact $\Delta$-space
has the needed countable Eberlein cover is itself an open problem reiterated
as Problem~5.6 in \cite{LS}.  Second, the ZFC rank-three non-$\Delta$ compactum
constructed in \cite[Corollary~4.6]{LS} did not yield an embedding into, or
representation as, a product of two compact $\Delta$-spaces.

Thus the source product problem remains open when neither factor has a PFNA
or a countable closed PFNA cover.  The result here is a substantial positive
subcase, not a full solution.

For novelty, the run's four lightweight indexes and bounded searches using
the arXiv id, title, authors, ``point-finite neighborhood assignment,''
``product,'' ``$\Delta$-space,'' and ``scattered Eberlein compact'' found
\cite{KL2,LS} but no explicit statement of Theorem~\ref{thm:pfna}.  This is a
candidate result awaiting expert review, not a definitive priority claim.

\section*{Human-review recommendation}

The core proof is short and should be checked first at the quantifier switch
from the fiberwise expansions to $W_n$.  The compact corollaries then require
only verification of the cited PFNA and countable-closed-union facts.  A
reviewer should also run a broader topology-database search for an earlier
abstract product lemma under the terminology ``scant cover.''

\begin{thebibliography}{9}
\bibitem{KL}
J. K\k{a}kol and A. Leiderman,
\emph{A characterization of $X$ for which spaces $C_p(X)$ are distinguished
and its applications}, Proc. Amer. Math. Soc. Ser. B \textbf{8} (2021),
86--99; arXiv:2011.14299.

\bibitem{KL2}
J. K\k{a}kol and A. Leiderman,
\emph{Basic properties of $X$ for which the space $C_p(X)$ is distinguished},
Proc. Amer. Math. Soc. Ser. B \textbf{8} (2021), 267--280;
arXiv:2104.10506.

\bibitem{LS}
A. Leiderman and P. Szeptycki,
\emph{On $\Delta$-spaces}, arXiv:2307.16047v1 (2023), accepted for
publication.
\end{thebibliography}

\end{document}
